% Separate post-CP36 candidate; independent of any unpromoted B100 assertion.
\subsection{Uniform wild level arcs and an exact infinite height-two family}

Let $p\ge5$, $q=p^2$, $N=q-1$. We retain the covariant
special-character convention of B62 and the manuscript
index convention that $\nu$ corresponds to actual
$p^{\nu+1}$-torsion.

\begin{lemma}[An integral arc with prescribed low and high steps]
\label{lem:B-RD101-arc}
\status{Candidate; complete proof awaiting independent audits.}
Let $p-1\mid e$, put $v=v_p(e)$, and suppose
\[
                n\ge v+2,\qquad e>p^{\,2n-3-v}.
                                                        \tag{B-RD101.1}
\]
There is a complete discretely valued mixed-characteristic
$(0,p)$ field $K$, with residue field
$\overline{\mathbf F}_p$, $\mu_p\subset K$ and absolute
index $e$, and a height-two BT group $G/K$ with a rational
point of exact order $p^n$.
It acquires good one-dimensional height-two formal reduction
over a cyclic tame extension $L/K$ of degree $N$.
Its special eigenvalues are $\zeta^{-A},\zeta^{-pA}$,
where $A>0$ is chosen with
\[
 \frac{e}{p^{\,2n-2-v}}<A<
       \frac{e}{p^{\,2n-3-v}},\qquad N\nmid A.
                                                        \tag{B-RD101.2}
\]
\end{lemma}
\begin{proof}
The required integer exists. Put
$x_0=e/p^{2n-2-v}$, so $px_0>1$.
If $x_0<1$, take $A=1$.
Otherwise take $\lfloor x_0\rfloor+1$, increasing by one
if this is divisible by $N$. This gives
$x_0<A\le x_0+2<px_0$, since $p\ge5$, and the
two successive choices cannot both be divisible by $N$.

Set
\[
 h=n-1-v\ge1,\quad U=q^hA,\quad
 b=\frac{eN}{p^v}-(p-1)U,\quad d=eN .
                                                        \tag{B-RD101.3}
\]
The integer $b$ is well-defined since $p^v\mid e$.
The two strict inequalities in (B-RD101.2) say exactly
\[
\begin{split}
 b-(q-p)q^{h-1}A
     &=\frac{N}{p^v}\bigl(e-p^{\,2h+v-1}A\bigr)>0,\\
 (q-p)q^hA-b
     &=\frac{N}{p^v}\bigl(p^{\,2h+v}A-e\bigr)>0 .
\end{split}                                             \tag{B-RD101.4}
\]
In particular $b>0$. These inequalities determine the
branches; no selection is made at an equality polygon.

Over $W=W(\overline{\mathbf F}_p)$ use the integral
height-two $p$-typical group $F_u$ of
Lemma~\ref{lem:B-RD62-arc}. Its exact logarithm and
coefficient bounds, from the original
\cite[equations (5.3)--(5.5), Proposition~5.7 and
(12.3)--(12.5)]{HopkinsGross1994Bundles}, are
\[
\begin{gathered}
 p l_j=l_{j-1}u^{p^{j-1}}+l_{j-2},\quad
 l_0=1,\quad l_{-1}=0,\\
 [p]_u(T)=pT+(1-p^{p-1})uT^p+\sum_{i>p}c_i(u)T^i,\\
 c_i(u)\in(p,u)\ (p<i<q),\qquad
 c_q(u)\equiv1\pmod{(p,u)} .
\end{gathered}                                           \tag{B-RD101.5}
\]
All coefficients are integral; those in degrees
$2,\ldots,p-1$ vanish. This is a full formal series,
not a polynomial approximation.
Define in $W[[X]]$
\[
 Y_0=X^A,\qquad Y_{j+1}=[p]_{X^b}(Y_j)
       \ (0\le j<n-1),\qquad
 P(X)=\Phi_{X^b}(Y_{n-1}),\quad\Phi_u(T)=[p]_u(T)/T .
                                                        \tag{B-RD101.6}
\]
Every composition is defined in the complete power-series ring.

Modulo $p$, the first inequality of (B-RD101.4) makes
the degree-$q$ term uniquely earliest for the first $h$
iterations, even at the largest input $q^{h-1}A$.
Thus their orders are $A,qA,\ldots,q^hA=U$.
At the next input $U$, if there is a next iteration,
the second inequality makes the degree-$p$ term
uniquely earliest. Each following input order increases,
so $b<(q-p)$ times that order remains strict.
The subsequent $v$ orders are therefore
\[
 J_r=p^rU+b\,\frac{p^r-1}{p-1}
                  \quad(0\le r\le v).                 \tag{B-RD101.7}
\]
For $v=0$ this simply says that the final order is $U$;
there are no degree-$p$ iterations. Intermediate terms
are always later than the chosen degree-$p$ term, because
their coefficients contain $u$ modulo $p$, and terms
of degree greater than $q$ are later than the degree-$q$
term. Each selected leading coefficient equals one
modulo $p$. It follows that every $Y_j$ has the
specified order with leading coefficient one.

The final input order is $J_v\ge U$ and
$b<(q-p)U\le(q-p)J_v$. Consequently the
$uY_{n-1}^{p-1}$ term is uniquely earliest in $P$, and
\[
 b+(p-1)J_v
   =p^v\{b+(p-1)U\}=eN=d .
 \qquad
 P(0)=p,\quad P(X)\equiv X^d+O(X^{d+1})\pmod p .
                                                        \tag{B-RD101.8}
\]
Weierstrass preparation gives a degree-$d$ Eisenstein
ring $R=W[[X]]/(P)$, with uniformizer $x$ and
$v_R(p)=d$. The chosen order sequence strictly increases,
and $d=b+(p-1)J_v>J_v$. Thus coefficients divisible by
$p$ cannot affect any iterate's leading value after
substituting $x$. In particular
$Y_{n-1}(x)\ne0$, while
$[p]^n(x^A)=Y_{n-1}(x)P(x)=0$.
This is an actual point of exact order $p^n$.

Take a primitive $N$th root $\zeta\in W^\times$ and
$a=\zeta^{-A}$. Since $q^h\equiv1\pmod N$ and
$e/p^v$ is an integer,
\[
                         b\equiv-(p-1)A\pmod N .
                                                        \tag{B-RD101.9}
\]
The exact logarithm recursion gives
$l_j(a^{p-1}u)=a^{p^j-1}l_j(u)$:
the second summand has the same weight because
$N\mid p^j-p^{j-2}$.
Hence $aT$ is an integral isomorphism
$F_{a^{p-1}u}\to F_u$ and
\[
 [p]_{a^{p-1}u}(a^{-1}T)=a^{-1}[p]_u(T),\qquad
 \Phi_{a^{p-1}u}(a^{-1}T)=\Phi_u(T).                  \tag{B-RD101.10}
\]
Together with (B-RD101.9), this proves
$Y_j(\zeta X)=a^{-1}Y_j(X)$ for every $j$ and the
exact identity $P(\zeta X)=P(X)$.
Thus $\sigma(x)=\zeta x$ defines an automorphism of
$R$ of order $N$. With $L=\operatorname{Frac}(R)$
and $K=L^{\langle\sigma\rangle}$, the actual extension
$L/K$ is tame of degree $N$, and $v_K(p)=d/N=e$.

The isomorphism $aT$ is a semilinear descent cocycle
whose $N$fold product is one. It fixes the marked point
because $a\,\sigma(x^A)=x^A$.
Effective descent of the generic compatible finite
levels gives $G/K$ with the stated rational point.
Its special Honda eigenvalues are $a,a^p$, with
no faithfulness assertion when $\gcd(A,N)>1$.
Finally $p-1\mid e$ and the algebraically closed residue
field permit a Hensel $(p-1)$st root of the unit in
$-p=\pi_K^e u$. The unique tame cyclotomic field
$\operatorname{Frac}(W)(\mu_p)$ is consequently
contained in $K$. All assertions follow.
\end{proof}

\begin{proposition}[Actual ambiguity at every allowed arc level]
\label{prop:B-RD101-ambiguity}
\status{Candidate; complete proof awaiting independent audits.}
Under (B-RD101.1), there exist two allowed height-two
BT groups $M,H$ over such a field $K$ with
$M[p^n]\simeq H[p^n]$ and different finite special
inertia characters. The first group individually acquires
supersingular good reduction; the second has a two-step
potentially good height-one filtration.
For every $h'\ge2$, $\nu_{\min}(p,e,h')\ge n$.
\end{proposition}
\begin{proof}
Let $P$ be the actual rational point of the lemma and
$C=\langle P\rangle\simeq\mathbf Z/p^n$.
Form the actual quotient $G'=G/C$, using the finite
flat closure in the good model over $\mathcal O_L$.
It remains potentially good. As proved in full in
Theorem~\ref{thm:B-RD62}, the determinant is
$\det T_pG=\mathbf Z_p(1)\otimes\psi$ with finite
$\psi$ of order dividing $p-1$, value
$\psi(\sigma)=a^{p+1}$, and trivial restriction to $G_L$.
The actual snake sequence
\[
 0\longrightarrow C\longrightarrow G[p^n]
 \longrightarrow G'[p^n]\longrightarrow C\longrightarrow0
                                                        \tag{B-RD101.11}
\]
and primitivity of its first constant line identify the
image as $(\mathbf Z/p^n)(\chi_p\psi)$.
Thus $M=G'\otimes\psi^{-1}$ has an actual truncation
\[
 0\longrightarrow\mu_{p^n}\longrightarrow M[p^n]
 \longrightarrow(\mathbf Z/p^n)(\psi^{-1})
 \longrightarrow0 .                                   \tag{B-RD101.12}
\]
The complete Kummer, prime-to-$p$ averaging and compatible
cocycle argument of Lemma~\ref{lem:B-RD62-extension}
lifts this exact extension to $H$, with graded groups
$\mu_{p^\infty}$ and
$(\mathbf Q_p/\mathbf Z_p)(\psi^{-1})$.
This uses no good-reduction claim for the whole $H$.

With $\lambda=a^{-1}$, the two special multisets are
$\{\lambda,\lambda^p\}$ and
$\{1,\lambda^{p+1}\}$. Their trace difference is
$-(\lambda-1)(\lambda^p-1)$, nonzero because
$N\nmid A$ and $\gcd(p,N)=1$.
Both groups satisfy exactly the given monodromy condition.
Adding the same constant group gives the claim for $h'>2$.
Failure at actual $p^n$ means failure at manuscript
index $n-1$, so the lower bound is $n$.
\end{proof}

\begin{corollary}[The exact existential mixed threshold]
\label{cor:B-RD101-threshold}
\status{Candidate; complete proof awaiting independent audits.}
Fix $p\ge5$, $p-1\mid e$, $v=v_p(e)$ and $n\ge v+2$.
There exists an allowed whole-supersingular/filtered
height-two comparison at actual level $p^n$ with different
finite special characters, across fields of absolute index
$e$, if and only if $e>p^{2n-3-v}$.
\end{corollary}
\begin{proof}
Necessity is Proposition~\ref{prop:B-RD76-mixed},
including its strict endpoint. Sufficiency is the
explicit preceding construction. No conclusion for
$n<v+2$ is being drawn from this statement.
\end{proof}

\begin{theorem}[An exact family at all wild valuations]
\label{thm:B-RD101-exact}
\status{Candidate; complete proof awaiting independent audits.}
For every prime $p\ge5$ and integer $v\ge0$, in the
full B.4 height-two setting with perfect residue field
and $\mu_p\subset K$,
\[
              \nu_{\min}(p,2p^v(p-1),2)=v+2 .
                                                        \tag{B-RD101.13}
\]
\end{theorem}
\begin{proof}
Put $e=2p^v(p-1)$; its $p$-adic valuation is exactly $v$.
At $n=v+2$, the necessary and sufficient mixed threshold
is $e>p^{v+1}$, which holds because $2(p-1)>p$.
The arc can be chosen uniformly as
\[
 A=1,\quad h=1,\quad b=(p-1)(p^2-2),\quad
 d=2p^v(p-1)N .
                                                        \tag{B-RD101.14}
\]
It has one degree-$q$ iteration followed by $v$
degree-$p$ iterations. The preceding proposition
therefore proves the lower bound $v+2$.

For the upper bound take actual level $n'=v+3$.
Lemma~\ref{lem:B-RD76-determinant} applies since
$n'\ge v+2$ and detects the total dimension and common
finite determinant. Passing to the completed maximal
unramified base preserves the index and inertia.
As in the full proof of Corollary~\ref{cor:B-RD76},
dimensions zero and two are determined by residual
Brauer lifting and Cartier duality, respectively.
Two filtered groups of total dimension one are
determined residually because $\chi_p\equiv1\pmod p$.
A mixed comparison would require
\[
               e>p^{\,2n'-3-v}=p^{v+3},
                                                        \tag{B-RD101.15}
\]
which contradicts $2p^v(p-1)<p^{v+3}$.
Thus only two whole potentially supersingular groups
with common determinant remain.

For their actual quotient target,
Proposition~\ref{prop:B-RD64-formal} uses
$a=eN$ and $q_0=1$,
$q_{j+1}=\min\{a+q_j,pq_j+p-1\}$.
For $0\le j\le v+1$,
$a\ge2(p-1)p^j$, since $N>p$.
Induction consequently gives
\[
 q_j=2p^j-1\quad(0\le j\le v+2),\qquad
 q_{v+2}-\frac a{p-1}=2p^v-1>0 .
                                                        \tag{B-RD101.16}
\]
This is exactly the required $q_{n'-1}$ threshold.
The actual-target weighted logarithm and graph proof
therefore recovers all finite special characters at
$p^{v+3}$, or manuscript index $v+2$.
Together with the explicit lower example this proves equality.
\end{proof}

\begin{remark}[Exact scope of the uniform extension]
\label{rem:B-RD101-scope}
\status{Candidate; no universal least-index formula is asserted.}
The construction covers the specified strict mixed-comparison
threshold when $n\ge v_p(e)+2$; it does not classify
two whole potentially good groups at arbitrary ramification.
The exact least-index family fixes $e=2p^v(p-1)$.
The cases $v=0$ and $v=1$ agree with B64 and the
separately written B100 candidate, but no unpromoted
B100 result is used as a dependency here.
The proof uses the complete canonical B62 construction,
B64 actual-target estimate and B76 wild necessary
bound/dimension detection.
The original Hopkins--Gross recurrence and coefficient
passages freshly read for B100 are reused with their
precise recorded scope; no additional primary reading
or image inspection is claimed.
A standalone exact supplement will check all branch
identities and the upper recurrence, without inferring
effectivity from a finite range.
\end{remark}
